\section{The algorithm and its running time}
\label{sec:composition}

We now combine the constructions.  For clarity, the following pseudocode
lists the feasibility steps; every successful call also returns the
embeddings and intermediate variable assignments needed by its inverse.
\begin{quote}\small\ttfamily
SolveSEFE($G_1,G_2$):\\
\quad normalize colour connectivity and split union blocks;\\
\quad apply common-cycle separation-pair preprocessing;\\
\quad compute the reference P-orders, treating markers as real edges;\\
\quad subdivide markers;\\
\quad for each nontrivial normalized factor $I$: \\
\qquad compute the common SPQR forest;\\
\qquad compute common blocks, union components, attachments, and lobes;\\
\qquad construct the doubled embedding-tree instance;\\
\qquad synchronize common sources and mark common-block subtrees;\\
\qquad use whole-sector constraints for exact common P-pole pairs;\\
\qquad otherwise add a hub, restrict exceptional reads, and compare faces;\\
\qquad solve the complete-sector instance; reject if infeasible;\\
\qquad decode the colour blocks and apply face-list repair;\\
\quad glue the normalized factors and undo the augmentations;\\
\quad output compatible combinatorial embeddings.
\end{quote}

\begin{proof}[Proof of Theorem~\ref{thm:main}]
By Section~\ref{sec:normalization} and the BKR separation-pair
preprocessing, the original input is equivalent to an AND family of
normalized instances of total size $O(N)$.  All transformations are
constructive in both directions.

For a normalized instance, Theorem~\ref{gen:compilation} gives a
complete-sector instance.  The forward witness map takes a SEFE to a
satisfying embedding and constraint assignment.  The inverse witness map
takes a satisfying assignment to planar colour-block embeddings satisfying
the local-order condition and nonempty face lists.  Theorem~\ref{rep:thm:repair}
then constructs compatible embeddings of the two colours and all relative
positions of common components.  Theorem~\ref{cse:main} decides the
complete-sector instance and constructs its witness in polynomial time.
The composition is therefore sound and complete.

For the time bound, consider a normalized factor of size $N_j$.  The
charging argument in Section~\ref{rep:algorithm} gives $O(N_j)$ relevant
pairs $(K,Z)$ and colour lobes.  The auxiliary graphs and source-labelled
embedding trees have $O(N_j^2)$ total size.  By the count in
Section~\ref{gen:size}, the compiled instance has size
$M=O((N_j+1)^3)$ and is constructed in $O((N_j+1)^4)$ time.
Theorem~\ref{cse:main} solves an instance of size $M$ in $O((M+1)^{20})$
time.  Consequently the total cost per factor is $O((N_j+1)^{60})$.

If the factor sizes are $N_1,\ldots,N_t$, their total is $O(N)$, and
$\sum_j(N_j+1)^{60}=O((N+1)^{60})$.  The normalization inverse and the
face-list reconstruction have smaller polynomial bounds.  All common
bridges and isolated common vertices are retained throughout, and no
relative component position is inferred from rotations alone.  This
proves the theorem.\qedhere
\end{proof}
